\section{Improvement Plan}

\subsection{Completion Audit of the Previous Plan}

\begin{longtable}{@{}L{0.14\linewidth}L{0.13\linewidth}L{0.33\linewidth}L{0.32\linewidth}@{}}
\toprule
Previous priority & Status & Evidence in the current code & Remaining gap \\
\midrule
\endhead
Zero-impedance bus-merge contract & partial & Canonical preprocessing merges zero-impedance buses and propagates \texttt{bus\_merge\_map}; PF APIs unproject merged \(v_m/\theta\) back to original bus cardinality. & \texttt{unproject\_bus\_vector()} still maps by \texttt{orig\_pos = ext\_bus - 1}, which is only safe under contiguous 1..N external bus numbering. The contract is still implicit, not machine-enforced. \\
Result unprojection policy across APIs & partial & \texttt{solve\_power\_flow()}, \texttt{solve\_handle()}, and \texttt{solve\_power\_flow\_fdpf()} now unproject PF voltages when merges exist. & Non-voltage artifacts (branch-oriented tables, OPF/DPF/ONR side outputs) still do not expose one unified projected-vs-original indexing policy with explicit mapping metadata for every artifact. \\
DCDC directional Jacobian consistency & partial & Parity equality equations use iterate-dependent bidirectional efficiency (\(p_{in}\ge0 \Rightarrow \eta\), \(p_{in}<0 \Rightarrow 1/\eta\)). & Parity Jacobian still picks derivative sign from reference \texttt{p\_ref\_mw} (warm-start sign), not from current iterate \(x\). This can mis-linearize exactly when \(p_{in}\) changes direction. \\
ONR stable-id regression semantics & mostly complete & ONR outputs now emit stable \texttt{ACBranch::index}; projection-equivalence tests verify disjoint/open-closed partition and radiality. & The suite still lacks a strict projected-vs-canonical branch-id set equality assertion on a case with deliberately nontrivial branch indices. \\
Visualization verification depth & complete & Network/geo/dashboard/profile outputs carry stable ids, DC losses, and loading fields; \texttt{test\_visualization} now compares branch loss/loading and carbon intensity across JSON, workbook, HTML dashboard payloads, and Plotly figure inputs. & Contract locked by the payload-level numeric consistency test. \\
Notebook and benchmark release gating & complete & \texttt{technical\_notebook\_pdf} is part of \texttt{.github/workflows/release-gate.yml}; the gate builds contract tests, runs a solver benchmark smoke, and emits \texttt{release\_gate\_summary.json}. & Contract locked by \texttt{test\_paper\_study\_contracts} and \texttt{scripts/compact\_benchmark\_summary.py}. \\
\bottomrule
\end{longtable}

\subsection{Native Branch-and-Cut Proof-Tail Direction}

The highest-value Native B\&C work is no longer root LP objective alignment or
threshold tuning.  The root-certified path is now close to the intended
regime: on the UC root-certificate family, the vendored HiGHS certificate can
close the gap with one LP and no native tree.  The remaining performance gap is
the non-root proof tail: Native can produce local proof information, but the
critical question is whether that information becomes a reusable certificate
that raises the live frontier lower bound.

The proof invariant to optimize is
\[
  z_{\mathrm{LB}} = \min_{N \in \mathcal{F}_{\mathrm{proof}}} z_N.
\]
Incumbent quality alone cannot certify optimality.  Every major B\&C change
should therefore answer one question: did it lift, prune, or refresh the proof
frontier so that this lower bound moves toward the incumbent?

\textbf{Primary chain.}  The next proof-tail implementation work should be
organized around the following closed loop:
\[
  \text{local proof artifact}
  \rightarrow \text{resolved reason-side frontier}
  \rightarrow \text{bound-lifting certificate}
  \rightarrow \text{queue lower-bound movement}.
\]
The Native code already has the main hooks for this chain:
\texttt{resolve\_dual\_proof\_target\_bound\_from\_local\_trail()},
\texttt{resolve\_dual\_proof\_conflict\_from\_local\_trail()},
\texttt{BoundLiftingCertificate},
\texttt{NodeQueue::propagate\_bound\_lifting\_certificates()}, and
\texttt{NodeQueue::propagate\_by\_dual\_proof\_domain()}.  The remaining work is
to make the chain auditable and to verify that it moves the frontier rather
than merely increasing artifact counts.

\begin{longtable}{@{}L{0.18\linewidth}L{0.34\linewidth}L{0.38\linewidth}@{}}
\toprule
Area & Current implementation anchor & Next improvement target \\
\midrule
\endhead
Local trail resolution & Native already resolves target bounds and conflicts
through the local domain trail, including previous-bound relaxation and
activity revalidation fields. & Add explicit counters for target/conflict
attempts, successes, missing reasons, scope blocks, and activity failures.
Reject every artifact whose resolved frontier does not revalidate the proof
activity after substitution. \\
Bound-lifting certificates & \texttt{BoundLiftingCertificate} is a first-class
object with frontier literals, target bound, source literal, hash, depth, and
proof-activity audit values. & Strengthen the type contract so cross-frontier
consumption requires resolved-frontier, activity-revalidated, target-validated,
and space-validated flags.  Keep local artifacts and published certificates
separate in names and call paths. \\
Queue certificate replay & Queue replay can tighten domains, append
\texttt{DomainReasonBound}, prune infeasible nodes, mark LP refresh, and publish
movement statistics. & Treat replay success as lower-bound movement, not as
scan count.  For every pass, record queue lower bound before and after,
frontier-domain movement, LP-refresh count, prunes, and refreshed-node bound
lift. \\
Queue index consistency & Sequential and threaded queues maintain separate
priority and bound views through \texttt{priority\_}, \texttt{bounds\_}, and
\texttt{dfs\_bounds\_}; the direct dual-proof replay path rekeys nodes when
\texttt{node.bound} changes. & Add an audit path that recomputes the minimum
bound from live nodes and compares it with the indexed queue lower bound after
certificate replay, post-update callbacks, LP refresh, and prune passes. \\
Dual proof rows & \texttt{build\_full\_dual\_proof\_row()} and
\texttt{build\_reduced\_cost\_mass\_dual\_proof\_row()} already check identity,
stationarity, activity range, and mapped cost errors. & Attach an LP-state
signature to each row: basis hash, row-dual hash, reduced-cost hash,
nonbasic-side hash, objective, and standard-form signature.  Reject proof-row
reuse on basis, side, or space mismatch. \\
Root and non-root proof semantics & \texttt{root\_gap\_certificate\_ready} and
\texttt{require\_tree\_exhaustion\_certificate} already separate root shortcut
logic from strict tree exhaustion. & Make the result proof mode explicit
(vendored root certificate, native root gap certificate, strict tree
exhaustion) so future maintenance cannot replace strict exhaustion with an
ordinary gap shortcut. \\
Parallel proof frontier & \texttt{WorkStealingPool} and \texttt{ActiveNodeBounds}
expose lower-bound views for worker deques and active nodes. & Centralize the
parallel proof lower bound as the minimum of shared queue, worker deques, and
active nodes, and log those components when proof-tail mode is active. \\
\bottomrule
\end{longtable}

\textbf{Acceptance metrics.}  A non-root proof-tail benchmark should be judged
by movement metrics, not by artifact volume.  The minimum useful diagnostic set
is:
\begin{itemize}
  \item \texttt{resolved\_target\_attempts}, \texttt{resolved\_target\_success},
    \texttt{resolved\_conflict\_attempts}, and
    \texttt{resolved\_conflict\_success};
  \item \texttt{resolution\_activity\_failed},
    \texttt{resolution\_missing\_reason}, and
    \texttt{resolution\_scope\_blocked};
  \item \texttt{bound\_lifting\_certificates\_published} and
    \texttt{bound\_lifting\_certificates\_rejected\_audit};
  \item \texttt{certificate\_queue\_tightenings},
    \texttt{certificate\_queue\_prunes},
    \texttt{certificate\_queue\_lower\_bound\_lift\_nodes},
    \texttt{certificate\_queue\_lp\_refresh\_marked}, and
    \texttt{certificate\_queue\_frontier\_bound\_move\_sum};
  \item queue audit records of the form
    \texttt{before}, \texttt{after}, \texttt{changed\_nodes},
    \texttt{tightened}, \texttt{pruned}, \texttt{computed\_min}, and
    \texttt{indexed\_min}.
\end{itemize}

The proof-tail target is to observe all of the following on a case where the
vendored root frontier is disabled or unavailable:
\[
  \texttt{resolved\_target\_success} > 0,
  \quad
  \texttt{bound\_lifting\_certificates\_published} > 0,
  \quad
  \texttt{certificate\_queue\_lower\_bound\_lift\_nodes} > 0,
  \quad
  \Delta z_{\mathrm{LB}} > 0.
\]

\textbf{Execution order for Native B\&C.}
\begin{enumerate}
  \item Harden local trail resolution and activity revalidation before adding
    new proof sources.
  \item Publish only resolved, audited \texttt{BoundLiftingCertificate} objects;
    keep proof artifacts local until resolution succeeds.
  \item Make queue replay prove real movement: tighten/prune/mark LP refresh,
    then refresh marked nodes and audit the lower-bound index.
  \item Add proof-row LP-state signatures after the certificate path is
    moving the frontier.
  \item Optimize parallel proof drain only after the single-frontier
    certificate path shows lower-bound movement.
\end{enumerate}

\subsection{Current Code-Quality Priorities}

\textbf{Priority 1: Remove hidden index assumptions in time-series and merged-network paths.}

\textbf{Why this is now first.} The new UC\(\rightarrow\)OPF\(\rightarrow\)PF pipeline has
good functional coverage, but still contains bus-position assumptions (for example
\texttt{bus-1} indexing when replaying OPF voltages) while bus-merge unprojection also
assumes contiguous external indices. These are correctness risks on real planning
cases with nontrivial bus numbering.

\textbf{Minimal action list.}
\begin{itemize}
  \item Introduce one shared utility for \texttt{bus\_index -> vector\_position} mapping and use it in time-series replay + unprojection code paths.
  \item Extend \texttt{BusMergeMap} with explicit original-position mapping, and use that instead of \texttt{ext\_bus-1}.
  \item Add regression cases with non-contiguous indices for PF, OPF replay, and merged-network unprojection.
\end{itemize}

\textbf{Priority 2: Make parity DCDC Jacobian consistent with the active operating direction.}

\textbf{Why this matters.} Equality constraints already use iterate-dependent
efficiency; Jacobian entries still use reference-sign logic. This creates local
model inconsistency near direction reversals and can degrade IPM/Newton steps.

\textbf{Minimal action list.}
\begin{itemize}
  \item Compute \(\partial g / \partial p_{dcdc}\) from current iterate \(x\) (or a smooth transition near zero).
  \item Add finite-difference tests that force sign reversals during solving.
  \item Add one economic dispatch case where bidirectional DCDC operation is active.
\end{itemize}

\textbf{Priority 3: Clarify UC initial-state semantics and startup-cost accounting.}

\textbf{Why this matters.} UC currently filters generators by \texttt{in\_service}
and derives \(u_{init}\) from the same flag. This conflates ``available in model''
with ``initially online'', and can hide startup-cost errors at \(t=0\).

\textbf{Minimal action list.}
\begin{itemize}
  \item Separate availability (\texttt{in\_service}) from initial commitment (for example \texttt{u\_init} field or UC input override).
  \item Allow initially-off but available units to participate in UC.
  \item Add unit tests for \(t=0\) startup behavior and cold-start transitions.
\end{itemize}

\textbf{Priority 4: Make cross-validation metrics physically complete for AC/DC systems.}

\textbf{Why this matters.} Current OPF-vs-PF loss comparison uses simplified OPF
supply-demand arithmetic (AC-side only) and does not fully replay load-shedding on
component-level load models. This can distort loss and mismatch diagnostics.

\textbf{Minimal action list.}
\begin{itemize}
  \item Compute OPF-side loss with the same accounting boundary as PF-side metrics (AC+DC loads, converters, controllable demand).
  \item Replay load shedding consistently onto the component tables used by PF residual assembly.
  \item Add tolerance-based AC/DC cross-validation tests on loss and power-balance metrics.
\end{itemize}

\textbf{Priority 5: Close remaining DC model I/O contract gaps.}

\textbf{Why this matters.} Runtime now uses \texttt{dc\_pv\_arrays} and
\texttt{dc\_static\_generators} paths, but top-level JSON I/O still serializes only
\texttt{dc.static\_generators}. This creates silent data loss when exchanging cases.

\textbf{Minimal action list.}
\begin{itemize}
  \item Add JSON (and Excel, if applicable) read/write support for \texttt{dc.pv\_arrays} and \texttt{dc.dc\_static\_generators}.
  \item Add round-trip tests that assert value equality for all DC generation subtypes.
  \item Update notebook field dictionary and consumer contract tables accordingly.
\end{itemize}

\textbf{Priority 6: Split benchmark-scale tests from fast correctness gates.}

\textbf{Why this matters.} \texttt{test\_unit\_commitment} is comprehensive but
benchmark-heavy; it is unsuitable as a fast regression gate and obscures real
correctness failures when CI timeouts occur.

\textbf{Minimal action list.}
\begin{itemize}
  \item Keep a fast UC smoke/contract test in default CI; move full benchmark comparison to nightly/perf jobs.
  \item Add explicit runtime budgets and labels for test tiers (smoke/contract/perf).
  \item Gate releases on \texttt{technical\_notebook\_pdf} + core contract tests.
\end{itemize}

\textbf{Priority 7: Add dedicated RPO regression tests.}

\textbf{Why this matters.}  The reactive-power optimisation (RPO) module is
fully implemented with a 4-phase outer-approximation MINLP, but has no
dedicated test file.  Convergence properties and OA gap behaviour are
currently exercised only through internal development runs.

\textbf{Minimal action list.}
\begin{itemize}
  \item Add at least one test case per RPO objective (loss minimisation, voltage deviation, combined).
  \item Verify that discrete tap/shunt positions satisfy integrality after convergence.
  \item Check that OA gap decreases monotonically across the 4 phases.
\end{itemize}

\textbf{Priority 8: Strengthen annual production simulation and lifecycle
validation.}

\textbf{Why this matters.}  Both the annual production simulation (L0/L2/L3
hierarchy) and the multi-year lifecycle simulation are implemented but only
have smoke-level test coverage.  There is no reference-case comparison
against an external production-cost tool, and lifecycle NPV results are not
cross-checked against manual calculations.

\textbf{Minimal action list.}
\begin{itemize}
  \item Add a small reference case (e.g.\ 5-bus, 24-hour scaled to 1~week) with known optimal cost from an external solver.
  \item Add lifecycle regression test comparing NPV against a spreadsheet on a known degradation curve.
  \item Verify that L0 energy budgets correctly constrain L2 weekly UC dispatch.
\end{itemize}

\subsection{Module-Level Investigation and Improvement Matrix}

The current project is now large enough that a single priority list is not
sufficient.  The table below records the module-by-module code-quality
investigation plan that should be updated whenever a major implementation path
changes.

\begin{longtable}{@{}L{0.17\linewidth}L{0.28\linewidth}L{0.45\linewidth}@{}}
\toprule
Module & Current investigation finding & Concrete improvement plan \\
\midrule
\endhead
Model / canonical projection & Rich public component model and useful canonical projection; remaining risk is hidden id/vector-position assumptions in replay, unprojection, and result consumers. & Introduce a shared id-position map utility; extend \texttt{BusMergeMap} with explicit original-position data; test non-contiguous bus and branch ids through PF, OPF, DPF, ONR, and exports. \\
I/O and result contracts & Input and output support is broad, but the runtime model now includes fields that are not uniformly serialized or reported. & Complete JSON/Excel support for DC PV arrays and DC static-generator subtypes; serialize parity OPF DER/storage/DCDC/flex vectors; add round-trip and golden-file tests. \\
Power flow & Main kernels are recently cleaned up: shared susceptance builder, split AC kernel, split Jacobian pattern, shared DC injection assembly.  Remaining risk is formula drift across PF, OPF, and validation code. & Make every converter and DC injection formula available through shared helpers; add finite-difference derivative tests for VSC and DCDC controls; keep FDPF and Newton susceptance/admittance construction under common builders. \\
Optimal power flow & Parity path is the strongest AC/DC NLP route; native generic path still has globalization/performance gaps; RPO has no dedicated regression file. & Fix DCDC directional derivative consistency around sign changes; add RPO loss/voltage/combined objective tests; keep parity-vs-native benchmark table current; expand output dictionaries. \\
Engine solvers & LP/MILP stack is mature; generic NLP IPM is improving but not yet parity-competitive for large OPF. & Separate smoke, contract, and performance test tiers; add persistent benchmark baselines for dual simplex, LP IPM, LCQP, branch-and-cut, and NLP IPM; keep presolve/postsolve round-trip checks mandatory. \\
Distribution PF & BFS and radial tests are credible, but coverage is still too feeder-simple for the full hybrid component catalog. & Add cross-validation feeders with transformers, switches, chargers, flexible loads, storage, asymmetric load, PV curves, and projected bus/branch ids. \\
Short circuit & IEC-style Zbus/fault reference coverage now includes a single-source named case for three-phase, single-line-ground, line-line, and line-line-ground faults, with explicit external-grid and zero component-bucket assertions. & Extend the named-reference set to transformer/vector-group and converter-rich feeders; keep contribution-bucket conservation checks mandatory for every new short-circuit component model. \\
Network reconfiguration & Stable branch ids and PF post-validation are mostly complete. & Add projected-vs-canonical branch-id set equality tests; test nontrivial physical branch ids; report infeasibility reasons consistently. \\
Carbon analysis & Static tracing, matrix solve, branch/VSC/DCDC loss attribution, and time-coupled ESS SOC/carbon-inventory dynamics are now covered by hand-calculated reference tests, including per-step system dispatch in dynamic series. & Add larger multi-storage chronological cases with nonzero self-discharge and mixed AC/DC storage, and keep explicit loss-bucket reference values whenever converter formulas change. \\
Time series / annual sim / lifecycle & End-to-end smoke tests are strong for regressions but not yet enough for benchmark-grade numerical claims. & Add a small external reference production-cost case; add spreadsheet-verified lifecycle NPV/degradation case; test UC initial commitment and startup accounting explicitly. \\
Reliability / resilience / MESS & Heuristic restoration is now pinned against a strict HiGHS-backed restoration MIP on a toy MESS feeder, and MESS movement/energy edge cases cover in-transit dispatch blocking, travel-energy infeasibility, and SOC clamping. & Broaden the strict reference tier to multi-island repair sequencing and rolling-horizon feeders; track heuristic-vs-MIP optimality gaps as benchmark metadata instead of ad hoc demo output. \\
Planning / market / stochastic & Paper-style studies are indexed by \texttt{reports/study\_contracts/paper\_studies.json}, with one reproducible CLI command, exact input manifest, reference result artifact, and smoke/reference regression per study. & Keep the manifest updated whenever a study command, input fixture, artifact schema, or regression target changes. \\
Visualization / reporting & JSON result contracts, workbooks, HTML dashboard payloads, and Plotly figure inputs now carry stable ids plus shared loss/loading/carbon metrics. & Keep \texttt{test\_visualization} as the consumer contract for new reporting fields and chart payloads. \\
Documentation / CI & Release branches are gated on notebook compilation, core contract tests, and compact benchmark summarization through \texttt{.github/workflows/release-gate.yml}. & Require docs updates when public formulas, result fields, or paper-study commands change. \\
\bottomrule
\end{longtable}

\subsection{Recommended Execution Order}

The highest-value technical sequence is now:
\begin{enumerate}
  \item remove index assumptions first (time-series replay + merge unprojection),
  \item fix parity DCDC Jacobian directional consistency and lock it with FD tests,
  \item solidify UC initial-state/startup semantics,
  \item align OPF-vs-PF cross-validation boundaries for AC/DC metrics,
  \item complete DC JSON/Excel contracts and corresponding notebook tables,
  \item split benchmark-scale tests from fast correctness gates,
  \item add dedicated RPO regression tests (convergence + integrality + OA gap),
  \item then strengthen annual production simulation and lifecycle validation with reference-case comparisons.
\end{enumerate}
